\chapter*{Conclusion}
\addcontentsline{toc}{chapter}{Conclusion}


This thesis studied the problem of fine-grained hierarchical bird image classification on the CUB-200-2011 dataset. Unlike flat species classification, hierarchical bird recognition requires the model to predict multiple taxonomy levels, including order, family, and species. This setting is challenging because bird species from the same family often share very similar global appearance and differ only in subtle local traits such as bill shape, wing pattern, head pattern, breast color, or tail shape. Therefore, an effective model should capture both hierarchical taxonomic structure and fine-grained visual evidence.

To address this problem, this thesis proposed an attribute-guided extension of H-CAST. The proposed method keeps H-CAST as the image-based hierarchical backbone and adds an attribute-guided learning module at the species-level branch. Instead of predicting all CUB attributes as a single flat vector, the model uses group-specific attribute classifiers based on semantic attribute groups. Each active attribute group is represented by a learnable query, which attends to species-level graph features produced by H-CAST. The resulting attribute-oriented features are used for auxiliary attribute prediction and are fused with the species class token for final species classification. Family and order predictions remain produced by the original H-CAST hierarchical stages.

A key component of the proposed method is the Hard-Sibling Discriminative Attribute Selection strategy. This strategy selects attributes that are useful for distinguishing species from visually similar sibling species within the same family. The goal is to reduce potentially redundant attribute supervision and focus the attribute branch on fine-grained cues that are more relevant to species-level discrimination. In addition, a diversity regularization term is used to encourage attribute-related attention to be less redundant.

Experiments were conducted under the official CUB train/test protocol. Compared with the H-CAST rerun, the proposed attribute-guided H-CAST improved Full-Path Accuracy from 82.86\% to 84.36\%, species accuracy from 84.55\% to 85.83\%, family accuracy from 95.22\% to 95.62\%, order accuracy from 98.74\% to 99.03\%, and wAP from 86.90\% to 88.00\%. TICE was also reduced from 4.35\% to 3.62\%. These results show that the proposed attribute-guided learning strategy provides useful additional supervision for hierarchical fine-grained bird classification.

The ablation studies further showed that the best FPA, species accuracy, and wAP among the fusion-design settings are obtained when attribute supervision, species-level fusion, and diversity regularization are used together. Attribute supervision alone, without species-level fusion, was not sufficient to improve species classification. Conversely, fusion without attribute supervision improved classification to some extent, but the attribute branch was not semantically aligned with CUB attributes. This confirms that both semantic attribute supervision and explicit fusion are important for the proposed method. The diversity loss improved species accuracy and wAP at a moderate weight, although it introduced a trade-off with taxonomy consistency.

The analysis of attribute selection also showed that the proposed HSDAS-199 setting achieved the best species accuracy and wAP among the evaluated attribute-selection variants. However, random attribute subsets remained competitive in some metrics, especially TICE and family accuracy. Therefore, the proposed attribute selection strategy should not be interpreted as dominating all alternatives in every metric. Instead, the results support HSDAS as a targeted and interpretable strategy that provides useful bias toward species-level discrimination.

The qualitative and diagnostic analysis provided additional insight into the model behavior. The proposed model corrected 236 species-level errors made by the H-CAST rerun, while introducing 161 new species-level errors where H-CAST was originally correct. This gives a net gain of 75 correctly classified test images, which is consistent with the improvement in species accuracy. Among the corrected cases, 194 were intra-family confusions, suggesting that many improvements occur in fine-grained cases between visually related species. At the same time, attention visualizations showed that the model should not be interpreted as a precise attribute localization model. Global attribute groups such as size and shape often received high pooling weights, and their spatial maps could be diffuse or include background regions. Thus, the attention maps were treated only as diagnostic evidence rather than ground-truth localization.

Despite the improvements, the proposed method has several limitations. First, it depends on the availability and quality of class-level attribute annotations. Second, the attribute-guided fusion is mainly applied to the species-level branch, while family and order predictions are not directly fused with attribute-oriented features. Third, the method does not use pixel-level attribute localization supervision, so the attention maps may reflect global morphology, silhouette, context, or background cues in addition to object-part evidence. Finally, the performance depends on the selected attribute subset and loss weights.

Future work can extend this study in several directions. One direction is to incorporate object masks, part annotations, or localization constraints to encourage attribute queries to focus more consistently on bird regions. Another direction is to improve taxonomy consistency by applying attribute-guided fusion not only at the species level but also at higher hierarchy levels. Further work can also investigate more robust attribute selection strategies, evaluate the method on additional fine-grained hierarchical datasets, and study whether attribute-guided representations improve interpretability under stronger localization supervision.

In summary, this thesis demonstrates that incorporating semantic attribute guidance into an H-CAST-based hierarchical classifier can improve fine-grained bird species recognition. The proposed method is most effective when attribute supervision is combined with species-level fusion and moderate diversity regularization. While the method does not provide reliable attribute localization, it improves hierarchical classification performance and offers a useful direction for combining visual grouping, taxonomy structure, and semantic attribute supervision in fine-grained recognition.